\PassOptionsToPackage{unicode=true}{hyperref} % options for packages loaded elsewhere
\PassOptionsToPackage{hyphens}{url}
\PassOptionsToPackage{dvipsnames,svgnames*,x11names*}{xcolor}
%
\documentclass[ignorenonframetext,]{beamer}
\setbeamertemplate{caption}[numbered]
\setbeamertemplate{caption label separator}{: }
\setbeamercolor{caption name}{fg=normal text.fg}
\beamertemplatenavigationsymbolsempty
\usepackage{lmodern}
\usepackage{amssymb,amsmath}
\usepackage{ifxetex,ifluatex}
\usepackage{fixltx2e} % provides \textsubscript
\ifnum 0\ifxetex 1\fi\ifluatex 1\fi=0 % if pdftex
  \usepackage[T1]{fontenc}
  \usepackage[utf8]{inputenc}
  \usepackage{textcomp} % provides euro and other symbols
\else % if luatex or xelatex
  \usepackage{unicode-math}
  \defaultfontfeatures{Ligatures=TeX,Scale=MatchLowercase}
\fi
% use upquote if available, for straight quotes in verbatim environments
\IfFileExists{upquote.sty}{\usepackage{upquote}}{}
% use microtype if available
\IfFileExists{microtype.sty}{%
\usepackage[]{microtype}
\UseMicrotypeSet[protrusion]{basicmath} % disable protrusion for tt fonts
}{}
\IfFileExists{parskip.sty}{%
\usepackage{parskip}
}{% else
\setlength{\parindent}{0pt}
\setlength{\parskip}{6pt plus 2pt minus 1pt}
}
\usepackage{xcolor}
\usepackage{hyperref}
\hypersetup{
            pdftitle={Power and sample size calculations},
            pdfauthor={Dave Harrington},
            colorlinks=true,
            linkcolor=darkblue,
            citecolor=darkblue,
            urlcolor=darkblue,
            breaklinks=true}
\urlstyle{same}  % don't use monospace font for urls
\newif\ifbibliography
\usepackage{color}
\usepackage{fancyvrb}
\newcommand{\VerbBar}{|}
\newcommand{\VERB}{\Verb[commandchars=\\\{\}]}
\DefineVerbatimEnvironment{Highlighting}{Verbatim}{commandchars=\\\{\}}
% Add ',fontsize=\small' for more characters per line
\usepackage{framed}
\definecolor{shadecolor}{RGB}{248,248,248}
\newenvironment{Shaded}{\begin{snugshade}}{\end{snugshade}}
\newcommand{\AlertTok}[1]{\textcolor[rgb]{0.94,0.16,0.16}{#1}}
\newcommand{\AnnotationTok}[1]{\textcolor[rgb]{0.56,0.35,0.01}{\textbf{\textit{#1}}}}
\newcommand{\AttributeTok}[1]{\textcolor[rgb]{0.77,0.63,0.00}{#1}}
\newcommand{\BaseNTok}[1]{\textcolor[rgb]{0.00,0.00,0.81}{#1}}
\newcommand{\BuiltInTok}[1]{#1}
\newcommand{\CharTok}[1]{\textcolor[rgb]{0.31,0.60,0.02}{#1}}
\newcommand{\CommentTok}[1]{\textcolor[rgb]{0.56,0.35,0.01}{\textit{#1}}}
\newcommand{\CommentVarTok}[1]{\textcolor[rgb]{0.56,0.35,0.01}{\textbf{\textit{#1}}}}
\newcommand{\ConstantTok}[1]{\textcolor[rgb]{0.00,0.00,0.00}{#1}}
\newcommand{\ControlFlowTok}[1]{\textcolor[rgb]{0.13,0.29,0.53}{\textbf{#1}}}
\newcommand{\DataTypeTok}[1]{\textcolor[rgb]{0.13,0.29,0.53}{#1}}
\newcommand{\DecValTok}[1]{\textcolor[rgb]{0.00,0.00,0.81}{#1}}
\newcommand{\DocumentationTok}[1]{\textcolor[rgb]{0.56,0.35,0.01}{\textbf{\textit{#1}}}}
\newcommand{\ErrorTok}[1]{\textcolor[rgb]{0.64,0.00,0.00}{\textbf{#1}}}
\newcommand{\ExtensionTok}[1]{#1}
\newcommand{\FloatTok}[1]{\textcolor[rgb]{0.00,0.00,0.81}{#1}}
\newcommand{\FunctionTok}[1]{\textcolor[rgb]{0.00,0.00,0.00}{#1}}
\newcommand{\ImportTok}[1]{#1}
\newcommand{\InformationTok}[1]{\textcolor[rgb]{0.56,0.35,0.01}{\textbf{\textit{#1}}}}
\newcommand{\KeywordTok}[1]{\textcolor[rgb]{0.13,0.29,0.53}{\textbf{#1}}}
\newcommand{\NormalTok}[1]{#1}
\newcommand{\OperatorTok}[1]{\textcolor[rgb]{0.81,0.36,0.00}{\textbf{#1}}}
\newcommand{\OtherTok}[1]{\textcolor[rgb]{0.56,0.35,0.01}{#1}}
\newcommand{\PreprocessorTok}[1]{\textcolor[rgb]{0.56,0.35,0.01}{\textit{#1}}}
\newcommand{\RegionMarkerTok}[1]{#1}
\newcommand{\SpecialCharTok}[1]{\textcolor[rgb]{0.00,0.00,0.00}{#1}}
\newcommand{\SpecialStringTok}[1]{\textcolor[rgb]{0.31,0.60,0.02}{#1}}
\newcommand{\StringTok}[1]{\textcolor[rgb]{0.31,0.60,0.02}{#1}}
\newcommand{\VariableTok}[1]{\textcolor[rgb]{0.00,0.00,0.00}{#1}}
\newcommand{\VerbatimStringTok}[1]{\textcolor[rgb]{0.31,0.60,0.02}{#1}}
\newcommand{\WarningTok}[1]{\textcolor[rgb]{0.56,0.35,0.01}{\textbf{\textit{#1}}}}
\usepackage{longtable,booktabs}
\usepackage{caption}
% These lines are needed to make table captions work with longtable:
\makeatletter
\def\fnum@table{\tablename~\thetable}
\makeatother
\usepackage{graphicx,grffile}
\makeatletter
\def\maxwidth{\ifdim\Gin@nat@width>\linewidth\linewidth\else\Gin@nat@width\fi}
\def\maxheight{\ifdim\Gin@nat@height>\textheight\textheight\else\Gin@nat@height\fi}
\makeatother
% Scale images if necessary, so that they will not overflow the page
% margins by default, and it is still possible to overwrite the defaults
% using explicit options in \includegraphics[width, height, ...]{}
\setkeys{Gin}{width=\maxwidth,height=\maxheight,keepaspectratio}
% Prevent slide breaks in the middle of a paragraph:
\widowpenalties 1 10000
\raggedbottom
\setbeamertemplate{part page}{
\centering
\begin{beamercolorbox}[sep=16pt,center]{part title}
  \usebeamerfont{part title}\insertpart\par
\end{beamercolorbox}
}
\setbeamertemplate{section page}{
\centering
\begin{beamercolorbox}[sep=12pt,center]{part title}
  \usebeamerfont{section title}\insertsection\par
\end{beamercolorbox}
}
\setbeamertemplate{subsection page}{
\centering
\begin{beamercolorbox}[sep=8pt,center]{part title}
  \usebeamerfont{subsection title}\insertsubsection\par
\end{beamercolorbox}
}
\AtBeginPart{
  \frame{\partpage}
}
\AtBeginSection{
  \ifbibliography
  \else
    \frame{\sectionpage}
  \fi
}
\AtBeginSubsection{
  \frame{\subsectionpage}
}
\setlength{\emergencystretch}{3em}  % prevent overfull lines
\providecommand{\tightlist}{%
  \setlength{\itemsep}{0pt}\setlength{\parskip}{0pt}}
\setcounter{secnumdepth}{0}

% set default figure placement to htbp
\makeatletter
\def\fps@figure{htbp}
\makeatother


\usepackage{amsmath,verbatim}

\usepackage{fancyvrb}
\usepackage{manfnt}
\usepackage[normalem]{ulem}

%\usepackage[colorlinks=true]{hyperref}

\mode<presentation>{\usetheme{Malmoe}}

%\synctex=1

\setbeamertemplate{headline}{}


\setbeamerfont{footline}{size=\scriptsize}
\setbeamerfont{frametitle}{shape=\scshape}
\setbeamertemplate{itemize items}[circle]
\setbeamercovered{transparent}

\setbeamertemplate{navigation symbols}{}
\setbeamertemplate{footline}[frame number]{} 


\definecolor{forest}{rgb}{0, .5, 0}
\definecolor{brick}{rgb}{.5, 0, 0}
\definecolor{darkgreen}{rgb}{0, .5, 0}
\definecolor{darkred}{rgb}{.7, .15, .15}
\definecolor{darkblue}{rgb}{0, 0, .5}
\definecolor{Green}{rgb}{0.2,1,0.2}


\newcommand{\R}{\textsf{R}}
\newcommand{\RStudio}{\textsl{R Studio}}
\newcommand{\var}{\mbox{Var}}
\newcommand{\Cov}{\mbox{Cov}}
\newcommand{\Cor}{\mbox{Cor}}
\newcommand{\cip}{\mbox{$\perp\!\!\!\perp$}}
\newcommand{\bpsi}{\mbox{\boldmath$\psi$}}
\newcommand{\bbeta}{\mbox{\boldmath$\beta$}}
\newcommand{\bhat}{\mbox{\boldmath$\hat\beta$}}
\newcommand{\btheta}{\mbox{\boldmath$\theta$}}
\newcommand{\beeta}{\mbox{\boldmath$\eta$}}
\newcommand{\bfeta}{\mbox{\boldmath$\eta$}}
\newcommand{\bpsinot}{\mbox{\boldmath{$\psi_0$}}}
\newcommand{\bZ}{\bf Z}

%\usepackage[english]{babel}
\usepackage{lmodern}
\renewcommand\ttfamily{\usefont{T1}{lmtt}{m}{n}}

%\usepackage{palatino}
\usepackage[T1]{fontenc}
%\usepackage[lmodern]{babel}

% make all tt fonts bold to look more like Verbatim
\renewcommand\ttfamily{\usefont{T1}{lmtt}{m}{n}}

% Comment these out if you don't want a slide with just the
% part/section/subsection/subsubsection title:
\AtBeginPart{
  \let\insertpartnumber\relax
  \let\partname\relax
  \frame{\partpage}
}
\AtBeginSection{
  \let\insertsectionnumber\relax
  \let\sectionname\relax
  \frame{\sectionpage}
}
\AtBeginSubsection{
  \let\insertsubsectionnumber\relax
  \let\subsectionname\relax
  \frame{\subsectionpage}
}

\begin{comment}

%reduces space between code chunk and output
\usepackage{etoolbox}
\setlength{\parskip}{0pt}
\setlength{\OuterFrameSep}{0pt}
\makeatletter
\preto{\@verbatim}{\topsep=-10pt \partopsep=10pt }
\makeatother

\end{comment}

\title{Power and sample size calculations}
\author{Dave Harrington}
\date{May 14 - 18, 2018}

\begin{document}
\frame{\titlepage}

\begin{frame}
\tableofcontents[hideallsubsections]
\end{frame}
\hypertarget{background}{%
\section{Background}\label{background}}

\begin{frame}{%
\protect\hypertarget{an-example}{%
An example}}

\includegraphics[width=0.7\textwidth,height=\textheight]{../figures/sprint.pdf}

See publication by
\href{../../clinical_papers/SPRINT_nejm.pdf}{\textcolor{darkblue}{SPRINT writing committee}}.

\end{frame}

\begin{frame}{%
\protect\hypertarget{design-of-sprint}{%
Design of SPRINT}}

From the methods section of the paper:

\vspace{2em}

\begin{quote}
We planned a 2-year recruitment period, with a maximum follow-up of 6 years, and anticipated a loss to follow-up of 2\% per year. With an enrollment target of 9250 participants, we estimated  that the trial would have 88.7\% power to detect  a 20\% effect with respect to the primary outcome, assuming an event rate of 2.2\% in the standard-treatment group.
\end{quote}

\end{frame}

\begin{frame}{%
\protect\hypertarget{type-i-error}{%
Type I error}}

Type I error (alpha error)

\begin{itemize}
\item
  Probability that trial will report a false positive, i.e., claim a
  significant result when there is no treatment effect
\item
  Typically set no larger than 5\%
\item
  Depends on method of analysis, does not depend on sample size
\end{itemize}

\end{frame}

\begin{frame}{%
\protect\hypertarget{power}{%
Power}}

\begin{itemize}
\item
  Probability that the trial will report a true positive, i.e., claim a
  significant result when there is a treatment effect
\item
  Should be 80\% or greater
\item
  Depends on sample size, method of analysis, and size of treatment
  effect.
\item
  Power calculations relevant when study is designed
\item
  Power calculations have little value after a study is complete.

  \begin{itemize}
  \tightlist
  \item
    Precision measured through confidence intervals
  \end{itemize}
\end{itemize}

\end{frame}

\begin{frame}{%
\protect\hypertarget{overview-of-sample-size-for-censored-data}{%
Overview of sample size for censored data}}

This unit focuses on the power of tests based on the exponential
distribution and the log-rank test.

As in standard designs, the power depends on

\begin{itemize}
\item
  Type I error (significance level \(\alpha\))
\item
  Difference of interest, \(\Delta\), under an alternative hypothesis
  \(H_A\).
\end{itemize}

A notable difference from the usual scenario is that power depends on
the \textbf{number of failures} that will be observed, not the total
sample size.

\end{frame}

\begin{frame}{%
\protect\hypertarget{overview}{%
Overview \ldots}}

In practice, designing a survival study involves deciding how many
patients or individuals to enter, as well as how long they should be
followed.

Designs are usually either

\begin{itemize}
\item
  \emph{fixed sample size}, with the sample size determined in advance,
  or
\item
  \emph{sequential}, which incorporate the possibility of stopping early
  for efficacy or futility
\end{itemize}

Collett, Chapter 12 covers sample size calculations.

\end{frame}

\hypertarget{power-calculations-for-normally-distributed-observations}{%
\section{Power calculations for normally distributed
observations}\label{power-calculations-for-normally-distributed-observations}}

\begin{frame}{%
\protect\hypertarget{testing-for-differences-between-two-means}{%
Testing for differences between two means}}

Suppose data consist of:

\begin{itemize}
\tightlist
\item
  Group 1: \((Y_{11}, \dots Y_{1n_1})\)
\item
  Group 0: \((Y_{01}, \dots Y_{0n_0})\)
\end{itemize}

Assume \[
Y_{1j} \sim N(\mu_1,\sigma^2),  ~ Y_{0j} \sim N(\mu_0,\sigma^2)
\]

The usual objective is to test: \[ 
H_0:  \mu_1 = \mu_0 \Rightarrow  ~ H_0: \triangle = 0,~
\text{where } \triangle = \mu_1 - \mu_0
\]

\end{frame}

\begin{frame}{%
\protect\hypertarget{power-for-a-two-sample-normal}{%
Power for a two sample normal}}

The standard test is based on the \(Z\) statistic:

\[   Z = \frac{\overline{Y}_1 - \overline{Y}_0}
       {\sqrt{s^2 (\frac{1}{n_1} +\frac{1}{n_0}})} \] where \(s^2\) is
the pooled sample variance (assuming equal variances).

This test statistic has a \(N(0,1)\) distribution under \(H_0\). If the
sample sizes are equal in the two groups, \(n_0=n_1=n/2\), then:

\begin{align*}
Z &= \frac{\overline{Y}_1 - \overline{Y}_0}
{\sqrt{s^2 (\frac{1}{n/2} +\frac{1}{n/2}})} = 
\frac{\overline{Y}_1 - \overline{Y}_0}{2s/\sqrt{n}}
\end{align*}

\end{frame}

\begin{frame}{%
\protect\hypertarget{the-steps-for-calculating-sample-size}{%
The steps for calculating sample size}}

\begin{enumerate}
[1.]
\item
  Determine the critical value, \(c\), for rejecting the null when it is
  true.
\item
  Calculate the probability of rejecting the null when the alternative
  is true, substituting \(c\) from above.
\item
  Write the expression in terms of sample size for a given power.
\end{enumerate}

\end{frame}

\begin{frame}{%
\protect\hypertarget{step-1}{%
Step 1}}

Set the significance level, \(\alpha\), the probability of rejecting the
null hypothesis when it is true, and solve for \(c\): \begin{align*}
\alpha &=  P\left(|\overline{Y_1} - 
\overline{Y_0}| > c ~|~ H_0\right)\\[1ex]
&=  P\left(\frac{|\overline{Y_1} - \overline{Y_0}|}{2s/\sqrt{n}} > 
\frac{c}{2s/\sqrt{n}} ~|~ H_0 \right) \\[1ex]
&= P\left(|Z| > \frac{c}{2s/\sqrt{n}}\right) = 
2  \Phi\left(\frac{c}{2s/\sqrt{n}}\right) \\[1ex]
\text{so } z_{1-\alpha/2} &=  \frac{c}{2s/\sqrt{n}}\\[2ex]
\text{or  } c &= \frac{z_{1-\alpha/2} (2)(s)}{\sqrt{n}}
\end{align*}

Note that \(z_{\gamma}\) is the value such that
\(\Phi(z_{\gamma})=Pr(Z<z_{\gamma}) = \gamma\).

\end{frame}

\begin{frame}{%
\protect\hypertarget{step-2-power-as-a-function-of-sample-size}{%
Step 2: power as a function of sample size}}

Calculate the probability of rejecting the null when \(H_A\) is true.

Begin by writing down the probability of a Type II error: \begin{align*}
\beta &= P\left(\text{accept $H_0$} ~|~ H_A\right)\\
\text{so power } =  1-\beta &=   P\left(\text{reject $H_0$} ~|~ H_A\right)\\
&=P\left(|\overline{Y}_1 - \overline{Y}_0| > c ~|~ H_A\right)\\
&= P\left(\frac{|\overline{Y}_1 - \overline{Y}_0| - \Delta}{2s/\sqrt{n}} > 
\frac{c-\Delta}{2s/\sqrt{n}} ~|~ H_A\right)\\
&= P\left(Z > \frac{c-\Delta}{2s/\sqrt{n}}\right)\\
&= P\left(Z > z_{1-\alpha/2} - \frac{\Delta}{2s/\sqrt{n}}\right) \\
&= 1 - \Phi\left(z_{1-\alpha/2} - \frac{\Delta}{2s/\sqrt{n}}\right)
\end{align*}

\end{frame}

\begin{frame}{%
\protect\hypertarget{notes-on-power}{%
Notes on power \ldots}}

The power is an increasing function of the standardized difference:
\[ \mu_T(\triangle) =  \frac{\triangle}{2s / \sqrt{n}} \]

This is the number of standard errors between the two means, under the
assumption of equal variances.

\begin{enumerate}
\item As $n$ increases, the power increases.
\item For fixed $n$, the power increases with $\triangle$.
\item For fixed $n$ and $\triangle$, the power decreases with $s$.
\item Assigning equal numbers of patients to the two groups 
$(n_1=n_0=n/2)$ is best for maximizing power.
\end{enumerate}

\end{frame}

\begin{frame}{%
\protect\hypertarget{step-3-sample-size-as-a-function-of-power}{%
Step 3: sample size as a function of power}}

From the calculation for power:

\begin{align*}
1 - \beta &= 1 - \Phi\left(z_{1-\alpha/2} - \frac{\Delta}{2s/\sqrt{n}}\right) \\
&\Longrightarrow \beta = \Phi\left(z_{1-\alpha/2} - \frac{\Delta}{2s/\sqrt{n}}\right) \\
&\Longrightarrow z_{1 - \beta} = z_{1 - \alpha/2} - \frac{\Delta}{2s/\sqrt{n}} \\
&\Longrightarrow z_{\beta} + z_{1 - \alpha/2} = \frac{\Delta}{2s/\sqrt{n}} \\
&\Longrightarrow n  = \frac{(z_{1-\alpha/2} + z_{1-\beta})^2 4} {s^2} \\
\end{align*}

\end{frame}

\begin{frame}{%
\protect\hypertarget{notes-on-sample-size}{%
Notes on sample size}}

\begin{enumerate}
\item  Sample size increases as $s$ decreases.
\item  Sample size increases as power increases.
\item  Sample size increases as $\alpha$ decreases.

\end{enumerate}

\end{frame}

\begin{frame}{%
\protect\hypertarget{an-example-1}{%
An example}}

Derive the total sample size required for 90\% power for detecting a
difference of 0.5 standard deviations between means, based on a
two-sided 0.05 level test.

\begin{center}
\begin{tabular}{rcl}
$\alpha$                   & = & ~0.05 \\[1ex]
$z_{1-\frac{\alpha}{2}}$   & = & ~1.96 \\[1ex]
$\beta$                    & = & ~0.10 \\[1ex]
$z_{1 - \beta} = z_{0.90}$ & = & ~1.28 \\[1ex]
\end{tabular}
\end{center}

\begin{align*}
n &=  \frac{\left(1.96  + 1.28 \right)^2
4 s^2}{\triangle^2} ~ \approx ~ \frac{42 ~s^2}{\triangle^2}
\end{align*}

For a 0.5 standard deviation difference, \(\Delta/s=0.5\),
\begin{eqnarray*}
n & \approx & \frac{42}{(0.5)^2} = 168
\end{eqnarray*}

\end{frame}

\begin{frame}[fragile]{%
\protect\hypertarget{calculations-in-r}{%
Calculations in R}}

In practice, the test statistic is a \(t\)-statistic, not a
\(Z\)-statistic.

The difference is small when sample sizes are large, but more important
for smaller sample sizes.

Best to use software based on the \(t\)-distribution. \footnotesize

\begin{Shaded}
\begin{Highlighting}[]
\KeywordTok{power.t.test}\NormalTok{(}\DataTypeTok{delta =} \FloatTok{.5}\NormalTok{, }\DataTypeTok{sd =} \DecValTok{1}\NormalTok{, }\DataTypeTok{sig.level =} \FloatTok{0.05}\NormalTok{,}
             \DataTypeTok{power =} \FloatTok{0.9}\NormalTok{, }\DataTypeTok{type =} \StringTok{"two.sample"}\NormalTok{,}
             \DataTypeTok{alternative =} \StringTok{"two.sided"}\NormalTok{, }\DataTypeTok{strict =} \OtherTok{TRUE}\NormalTok{)}
\end{Highlighting}
\end{Shaded}

\begin{verbatim}
## 
##      Two-sample t test power calculation 
## 
##               n = 85.031
##           delta = 0.5
##              sd = 1
##       sig.level = 0.05
##           power = 0.9
##     alternative = two.sided
## 
## NOTE: n is number in *each* group
\end{verbatim}

\end{frame}

\hypertarget{censored-data}{%
\section{Censored data}\label{censored-data}}

\begin{frame}{%
\protect\hypertarget{sample-size-based-on-the-log-rank-test}{%
Sample size based on the log-rank test}}

\emph{Recap of the log-rank}

Consider a two-group survival problem, with equal numbers of individuals
in the two groups (\(n_0\) in group 0 and \(n_1\) in group 1).

Let \(\tau_1,...,\tau_K\) represent the \(K\) ordered, distinct failure
times, and at the \(j\)-th event time:

\begin{center}
\begin{tabular}{cccc}
\hline
& \multicolumn{2}{c}{Fail} & \\ \cline{2-3}
\multicolumn{1}{c}{Group } & ~~~Yes~~~ & ~~~No~~~ & Total\\  \hline
0 & $d_{0j}$ & $r_{0j} - d_{0j}$ & $r_{0j}$ \\[1ex]
1 & $d_{1j}$ & $r_{1j} - d_{1j}$ & $r_{1j}$ \\[1ex]
\hline
Total &  $d_j  $ & $r_j   - d_j  $ & $r_j  $  \\ 
 \hline
\end{tabular}
\end{center}

where \(d_{0j}\) and \(d_{1j}\) are the number of events in group 0 and
1, respectively, at the \(j\)-th event time, and \(r_{0j}\) and
\(r_{1j}\) are the corresponding numbers at risk.

\end{frame}

\begin{frame}{%
\protect\hypertarget{the-log-rank-test-statistic-z-statistic-version}{%
The log-rank test statistic (\(Z\)-statistic version)}}

\begin{align*}
Z_{LR} &=   \frac{\sum_{j=1}^{K} (d_{1j} - e_j)}{ 
 \sqrt{\sum_{j=1}^{K} v_j}} \qquad \text{ with } e_j = d_j \frac{r_{1j}}{r_j} \\[1ex]
v_j &= r_{1j} r_{0j} d_j(r_j-d_j)/[r_j^2(r_j-1)]
\end{align*}

\end{frame}

\begin{frame}{%
\protect\hypertarget{distribution-of-the-log-rank-statistic}{%
Distribution of the log-rank statistic}}

Suppose that the hazard rates in the two groups are \(\lambda_0(t)\) and
\(\lambda_1(t)\), with hazard ratio
\[  \theta= e^\beta = \frac{\lambda_1(t)}{\lambda_0(t)} \]

and suppose \(H_0: \beta=\log(\theta)=0\), which is equivalent to
\(H_0: \theta=1\).

It is possible to show that if there are no ties, and the observed
distribution is ``near’’ \(H_0\),

then

\begin{itemize}
\item
  \(E(d_{1j} - e_j | d_{1j}, d_{0j}, r_{1j}, r_{0j}) \approx \log(\theta)/4\)
\item
  \(v_j \approx 1/4\)
\end{itemize}

\end{frame}

\begin{frame}{%
\protect\hypertarget{distribution-of-the-log-rank-statistic-1}{%
Distribution of the log-rank statistic \ldots}}

At a value \(\log(\theta)\) under the alternative: \begin{align*} 
Z_{LR} &\approx  \frac{\sum_{j=1}^K \log(\theta)/4}{\sqrt{\sum_{j=1}^K 1/4}}\\[1ex]
&=\frac{d \log(\theta)/4}{\sqrt{d/4}}\\[1ex]
&= \frac{\sqrt{d} \log(\theta)}{2}\\[1ex]
\text{and } Z_{LR}  &\sim  N(\frac{\sqrt{d} \log(\theta)}{2},1) 
\end{align*}

\end{frame}

\begin{frame}{%
\protect\hypertarget{power-of-the-log-rank-test}{%
Power of the log-rank test}}

Using a similar argument to before, the power of a two-sided, level
\(\alpha\) logrank test is approximately:

\[ \text{Power} (\theta) \approx 
1 - \Phi\left[z_{1- \frac{\alpha}{2}} - \frac{\sqrt{d} \log(\theta)}{2} \right]  \]

Power depends on only \(d\) and \(\theta\).

\vspace{0.3in}

Possible to solve for required number of events to achieve a certain
power at a specified value of \(\theta\)\ldots

\end{frame}

\begin{frame}{%
\protect\hypertarget{power-of-the-log-rank-test-1}{%
Power of the log-rank test\ldots}}

For \(\text{Power}(\theta) = 1-\beta\), \(d\) must satisfy
\begin{align*}
1 - \beta &= 
1 - \Phi \left(z_{1- \frac{\alpha}{2}} - \frac{\sqrt{d} \log(\theta)}{2}\right)\\[1ex]
\Rightarrow  z_{\beta} &=  
z_{1- \frac{\alpha}{2}} - \frac{\sqrt{d} \log(\theta)}{2}\\[1ex]
\Rightarrow  d  &=  
\frac{4 \left(z_{1- \frac{\alpha}{2}} - z_{\beta}\right)^2}
{[\log(\theta)]^2}\\[1ex]
\text{or } d &=  
\frac{4 \left(z_{1- \frac{\alpha}{2}} + z_{1-\beta}\right)^2}{[\log(\theta)]^2} 
\end{align*}

\end{frame}

\begin{frame}{%
\protect\hypertarget{example}{%
Example}}

Suppose investigators are planning a 2-arm study, and want to detect a
hazard ratio of 1.5 with 90\% power at a 2-sided significance level of
\(\alpha=0.05\).

Required number of events: \begin{align*}
d &= \frac{4 \left(z_{1- \frac{\alpha}{2}} + z_{1-\beta}\right)^2}
{[\log(\theta)]^2}\\[1ex]
&= \frac{4 (1.96 + 1.282)^2}{[\log(1.5)]^2}\\[1ex]
&\approx \frac{42}{0.1644} = 256
\end{align*}

\end{frame}

\begin{frame}{%
\protect\hypertarget{events-required-for-various-hazard-ratios}{%
Events required for various hazard ratios}}

\begin{center}
\begin{tabular}{ccc}
Hazard & \multicolumn{2}{c}{\bf Power} \\ \cline{2-3}
Ratio & ~~80\%~~ & ~~90\%~~ \\ \hline
1.5 & 191 & 256 \\
2.0 & ~66 & ~88 \\
2.5 & ~38 & ~50 \\
3.0 & ~26 & ~35 \\
\hline
\end{tabular}
\end{center}

Most studies are designed to detect a hazard ratio of 1.5-2.0.

\#Practical considerations

\end{frame}

\begin{frame}{%
\protect\hypertarget{practical-considerations}{%
Practical considerations}}

\begin{itemize}
\tightlist
\item
  Deciding on \(\theta\)
\item
  Translating the number of failures to number of enrolled participants
\end{itemize}

Easiest to think about this for the hazard ratio \(\theta\) of two
exponential distributions.

If \(T_i \sim exp(\lambda_i)\), then
\[   \text{Median}(T_i) = -\log(0.5)/\lambda_i \]

It follows that \[
\frac{\text{Median}(T_1)}{\text{Median}(T_0)} = \frac{\lambda_0}{\lambda_1} = e^{-\beta} = \frac{1}{\theta}
\] Doubling the median survival of a treatment group compared to a
control group corresponds to halving the hazard.

\end{frame}

\begin{frame}{%
\protect\hypertarget{using-r-year-survival-rates-with-an-exponential-distribution}{%
Using \(R\)-year survival rates with an exponential distribution}}

Suppose the \(R\)-year survival rate is \(S_1(R)\) in group 1 and
\(S_0(R)\) in group 0.

Under the exponential model: \[ S_i(R) = \exp(-\lambda_i R) \]

Hence, \[   \frac{\log(S_1(R))}{\log(S_0(R))} 
= \frac{-\lambda_1 R}{-\lambda_0 R} =\frac{\lambda_1}{\lambda_0} =
e^{\beta} = \theta\]

Hence, doubling the hazard rate from group 1 to group 0 corresponds to
doubling the log of the \(R\)-year survival rate.

Note that this result does not depend on \(R\).

\end{frame}

\begin{frame}{%
\protect\hypertarget{example-1}{%
Example}}

Suppose the 5-year survival rate on treatment A is 20\% and
investigators want 90\% power to detect an improvement to 30\%.

The corresponding hazard ratio of treatment to control is:
\[ \frac{ln(0.3)}{ln(0.2)} = \frac{-1.204}{-1.609} = 0.748 \]

From the power formula for the log-rank, the number of events needed to
detect this improvement with 90\% power, based on a 2-sided 5\% level
test is \[    d = \frac{4 (1.96 + 1.282)^2}{[\log(0.748)]^2}  =  499 \]

\end{frame}

\begin{frame}{%
\protect\hypertarget{translating-to-number-of-enrolled-patients}{%
Translating to number of enrolled patients}}

Suppose a study enters all \(N\) patients at time 0, and will continue
the study for \(F\) units of time.

Under \(H_0\), the probability that an individual will fail during the
study is: \begin{align*}  
P(fail)   &=    \int_{0}^{F} \lambda_0 e^{-\lambda_0 t} dt\\[1ex]
   &=   1- e^{-\lambda_0 F}\\
\end{align*} Hence, if power calculations imply the study needs \(d\)
failures, then

\[  d = (N/2)(1 - e^{-\lambda_0 F}) + (N/2)(1 - e^{-\lambda_1 F})\]

\end{frame}

\begin{frame}{%
\protect\hypertarget{translating-to-number-of-enrolled}{%
Translating to number of enrolled\ldots}}

The solution for \(N\) requires values of \(F\) and \(d\).

\begin{enumerate}
[1.]
\item
  Assume a HR \(\theta\), then calculate \(d\).
\item
  Assume a follow-up period \(F\), then calculate \(N\).
\end{enumerate}

\end{frame}

\begin{frame}{%
\protect\hypertarget{example-2}{%
Example}}

Suppose investigators wish to detect a 50\% improvement in the median
survival from 12 months to 18 months with 80\% power at \(\alpha=0.05\),
and plan to follow participants for 3 years (36 months).

Use the two medians to calculate \(\lambda_0\), \(\lambda_1\), and the
hazard ratio, \(\theta\): \begin{align*}
\text{Median}(T_i) &= -\log(0.5)/\lambda_i \\[1ex]
\text{so } 
\lambda_1 &=  \frac{-\log(0.5)}{M1} = \frac{0.6931}{18} = 0.0385\\[1ex]
\lambda_0 &=  \frac{-\log(0.5)}{M0} = \frac{0.6931}{12} = 0.0578\\[1ex]
\theta &= \frac{\lambda_1}{\lambda_0} = \frac{0.0385}{0.0578} 
= \frac{12}{18} = 0.667
\end{align*}

From the earlier table, the number of events required is \(d=191\) (same
for \(\theta=1.5\) as it is for \(\theta = 1/1.5=0.667\)).

\end{frame}

\begin{frame}{%
\protect\hypertarget{example-3}{%
Example \ldots}}

Now solve\ldots \begin{align*}
191 &=  (N/2)(1 - e^{-0.0578(36)}) + (N/2)(1 - e^{-0.0385(36)})\\[1ex]
    &=  (N/2)(0.875) + (N/2)(0.7500) = (N/2)(1.625)\\[1ex]
\Rightarrow N &= 235
\end{align*}

Round up to 236 and randomize 118 patients to each treatment group.

\end{frame}

\begin{frame}{%
\protect\hypertarget{another-example}{%
Another example}}

Even if accrual does not all happen at time 0, this formula can be
surprisingly useful.

A clinical trial in esophageal cancer will randomize patients to
radiotherapy alone (Rx A) versus radiotherapy plus chemotherapy (Rx B).

\begin{itemize}
\item
  The goal of the study is to compare the two treatments with respect to
  survival, using the log-rank test.
\item
  From historical data, the median survival on Rx A for this disease is
  around 9 months.
\item
  Want 90\% power to detect an improvement in the median to 14 months.
\item
  Past studies have been able to accrue approximately 50 patients per
  year.
\end{itemize}

\end{frame}

\begin{frame}[fragile]{%
\protect\hypertarget{esophageal-cancer}{%
Esophageal cancer}}

\small

Start by estimating the number of events the trial will need for 90\%
power.

\footnotesize

\begin{Shaded}
\begin{Highlighting}[]
\CommentTok{# using years as time scale}
\NormalTok{alpha =}\StringTok{ }\FloatTok{0.05}\NormalTok{; beta =}\StringTok{ }\FloatTok{0.10}
\NormalTok{med}\FloatTok{.0}\NormalTok{ =}\StringTok{ }\FloatTok{0.75}\NormalTok{; med}\FloatTok{.1}\NormalTok{ =}\StringTok{ }\DecValTok{14}\OperatorTok{/}\DecValTok{12}

\NormalTok{num.events =}\StringTok{ }\DecValTok{4} \OperatorTok{*}\StringTok{ }\NormalTok{(}\KeywordTok{qnorm}\NormalTok{(}\DecValTok{1} \OperatorTok{-}\StringTok{ }\NormalTok{alpha}\OperatorTok{/}\DecValTok{2}\NormalTok{)}\OperatorTok{+}\StringTok{ }\KeywordTok{qnorm}\NormalTok{(}\DecValTok{1} \OperatorTok{-}\StringTok{ }\NormalTok{beta))}\OperatorTok{^}\DecValTok{2} \OperatorTok{/}
\StringTok{  }\KeywordTok{log}\NormalTok{(med}\FloatTok{.1}\OperatorTok{/}\NormalTok{med}\FloatTok{.0}\NormalTok{)}\OperatorTok{^}\DecValTok{2}
\NormalTok{num.events}
\end{Highlighting}
\end{Shaded}

\begin{verbatim}
## [1] 215.3
\end{verbatim}

\small

The trial has to enroll a minimum of 216 patients (with complete
follow-up).

If the anticipated accrual rate (50 patients per year) is correct, the
trial could enroll 250 patients over 4 years, then follow participants
until 216 events are observed.

\begin{itemize}
\tightlist
\item
  How long should participants be followed?
\end{itemize}

\end{frame}

\begin{frame}{%
\protect\hypertarget{a-more-realistic-accrual-pattern}{%
A more realistic accrual pattern}}

In reality, not everyone will enter the study on the same day.

Instead, the accrual will occur in a ``staggered’’ manner over a period
of time.

\emph{The standard assumption about enrollment}

Participants enter the study uniformly over an accrual period lasting
\(A\) units of time, and that after the accrual period, follow-up will
continue for another \(F\) units of time.

The translation \(d\) to \(N\) requires the probability that a
participant is observed to have an event under this accrual and
follow-up scenario.

\begin{align}
\nonumber
P(\text{fail}) &=  
\int_{0}^{A} P(\text{fail}|\text{enter at } a) f(a) da \\[1ex]
&= 1-\frac{\int_{0}^{A} S(a+F)~ da}{A}
\label{exactP}
\end{align}

\end{frame}

\begin{frame}{%
\protect\hypertarget{a-more-realistic-accrual-pattern-1}{%
A more realistic accrual pattern \ldots}}

Solve for \(d\), where \(P_c\) is the proportion of failures in the
control group and \(P_e\) is the proportion of failures in the
experimental group: \begin{align*}
d &= (N/2) P(\text{fail};\lambda_0) +(N/2) P(\text{fail};\lambda_1)\\
  &= (N/2) P_c + (N/2) P_e\\
  &= (N/2) (P_c + P_e)
\end{align*}

Solve for \(N\) based on the previous formula for \(d\): \begin{align*}
N &= \frac{2~d}{(P_c + P_e)}\\[1ex]
N &= \frac{8 ~ \left(z_{1-\frac{\alpha}{2}} + z_{1-\beta}\right)^2}
{[\log(\theta)]^2} \times \frac{1}{(P_c + P_e)}
\end{align*}

\end{frame}

\begin{frame}{%
\protect\hypertarget{calculating-p_c-and-p_e-from-equation}{%
Calculating \(P_c\) and \(P_e\) from equation (\ref{exactP})}}

If failure times are exponential distributed, (\ref{exactP}) implies:

\begin{align} \label{pExp}
P_i &= 1 - \frac{exp(-\lambda_i F) (1-exp(-\lambda_i A))}{\lambda_i A} \qquad \text{ for } i = c, e 
\end{align}

Freedman suggested an approximation for \(P_c\) and \(P_e\), by
computing the probability of an event at the median duration of
follow-up, \((A/2 + F)\): \begin{align*}
P_i = P(\text{fail};\lambda_i) &= 1 - exp[-\lambda_i (A/2 + F)]
\end{align*}

He showed that this approximation works pretty well for the exponential
distribution (i.e., it gives values close to (\ref{pExp})).

\end{frame}

\begin{frame}{%
\protect\hypertarget{other-approximations}{%
Other approximations}}

Rubenstein, Gail, and Santer (1981) and Lachin and Foulkes (1986) have
given more accurate approximations for calculating the probablilty of an
event.

\begin{itemize}
\item
  These methods are more complicated and require software
\item
  Software now widely available in \textsf{R} and other packages
\end{itemize}

\textsf{R} packages: \texttt{Hmisc}, \texttt{TrialSize},
\texttt{gsDesign}

\begin{itemize}
\item
  Many other \textsf{R} packages and programs
  \href{https://cran.r-project.org/web/views/ClinicalTrials.html}{here}
\item
  \texttt{cpower()} in \texttt{Hmisc} is a simple and robust program for
  computing power as a function of the ususal parameters
\item
  \texttt{nSurv()} more complicated to use, but allows many more options
\end{itemize}

Commercial packages: EaST, nQuery, SAS, etc.

\end{frame}

\begin{frame}{%
\protect\hypertarget{using-cpower}{%
Using cpower()}}

Parameters:

\begin{itemize}
\item
  \texttt{tref}: time point at which mortalities estimated, usually
  given in years
\item
  \texttt{n}: total sample size (both groups combined). If allocation is
  unequal so that there are not n/2 observations in each group, sample
  sizes can be specified in \texttt{nc} and \texttt{ni}.
\item
  \texttt{mc}: tref-year mortality, control group, as a decimal. This is
  the value of the control survivor function at \texttt{tref}.
\item
  \texttt{r}: relative \% reduction in mc by intervention. A reduction
  from 50\% to 40\% mortality at time tref is a 20\% reduction.
\item
  \texttt{accrual}: duration of accrual period, in same units as
  \texttt{tref}
\item
  \texttt{tmin}: minimum follow-up time
\end{itemize}

\end{frame}

\begin{frame}{%
\protect\hypertarget{using-cpower-1}{%
Using cpower()\ldots}}

\begin{itemize}
\item
  \texttt{noncomp.c}: \% non-compliant in control group (drop-ins)
\item
  \texttt{noncomp.i}: \% non-compliant in intervention group (drop-outs,
  non-adherers)
\item
  \texttt{alpha}: type I error probability. A 2-tailed test is assumed.
\item
  \texttt{nc}: number of subjects in control group
\item
  \texttt{ni}: number of subjects in intervention group. \texttt{nc} and
  \texttt{ni} are specified exclusive of n.
\item
  \texttt{pr}: set to FALSE to suppress printing of details
\end{itemize}

\end{frame}

\begin{frame}{%
\protect\hypertarget{using-cpower-for-esophageal-cancer-example}{%
Using cpower() for esophageal cancer example}}

The parameters:

\begin{itemize}
\tightlist
\item
  \texttt{tref} is 9 months, or 0.75 years, and since that is the median
  in the control group, \texttt{mc} = 0.5 = \(S_c(0.75)\)
\item
  \texttt{n} = 150 patients enrolled over 3 years
\item
  \texttt{accrual} = 3, and let us set \texttt{tmin} = 1 year of
  follow-up
\item
  \texttt{alpha} = 0.05\\
\item
  Assume no non-compliance
\end{itemize}

\end{frame}

\begin{frame}{%
\protect\hypertarget{using-cpower-for-esophageal-cancer-example-1}{%
Using cpower() for esophageal cancer example \ldots}}

Calculating \texttt{r}, using approximations \ldots

The median for the intervention is 14 months, so

\begin{itemize}
\item
  \(\lambda_i = \log(2)/(14/12) = 0.59\)
\item
  \(S_i(0.75) = \exp(-(0.59)(0.75)) = 0.64\)
\end{itemize}

Mortality at 0.75 years is reduced from 50\% to 36\%, a proportionate
reduction of (14/50) = 0.28, or 28\%.

\end{frame}

\begin{frame}[fragile]{%
\protect\hypertarget{using-cpower-for-esophageal-cancer-example-2}{%
Using cpower() for esophageal cancer example \ldots}}

What if we allow 1 year of follow-up?

\scriptsize

\begin{Shaded}
\begin{Highlighting}[]
\KeywordTok{library}\NormalTok{(Hmisc)}
\KeywordTok{cpower}\NormalTok{(}\DataTypeTok{tref =} \FloatTok{0.75}\NormalTok{, }\DataTypeTok{n =} \DecValTok{150}\NormalTok{, }\DataTypeTok{mc =} \FloatTok{0.5}\NormalTok{, }
       \DataTypeTok{r =} \DecValTok{28}\NormalTok{, }\DataTypeTok{accrual =} \DecValTok{3}\NormalTok{, }\DataTypeTok{tmin =} \DecValTok{1}\NormalTok{, }
       \DataTypeTok{noncomp.c =} \DecValTok{0}\NormalTok{, }\DataTypeTok{noncomp.i =} \DecValTok{0}\NormalTok{, }
       \DataTypeTok{alpha =} \FloatTok{0.05}\NormalTok{, }\DataTypeTok{pr =} \OtherTok{TRUE}\NormalTok{)   }
\end{Highlighting}
\end{Shaded}

\end{frame}

\begin{frame}[fragile]

\scriptsize

\begin{verbatim}
## 
## Accrual duration: 3 y  Minimum follow-up: 1 y
## 
## Total sample size: 150 
## 
## Alpha= 0.05 
## 
## 0.75-year Mortalities
##      Control Intervention 
##         0.50         0.36 
## 
## Hazard Rates
##      Control Intervention 
##      0.92420      0.59505 
## 
## Probabilities of an Event During Study
##      Control Intervention 
##      0.86581      0.74288 
## 
## Expected Number of Events
##      Control Intervention 
##         64.9         55.7 
## 
## Hazard ratio: 0.64386 
## Standard deviation of log hazard ratio: 0.18261
\end{verbatim}

\begin{verbatim}
##   Power 
## 0.67402
\end{verbatim}

\end{frame}

\begin{frame}{%
\protect\hypertarget{example-esophageal-cancer}{%
Example: Esophageal cancer}}

Power of initial guess is too low.

\begin{itemize}
\tightlist
\item
  67\% power because only 65 + 56 = 121 events
\end{itemize}

We can increase enrollment time or follow-up time or both.

\texttt{cpower()} can be used to search for a design.

\end{frame}

\begin{frame}[fragile]{%
\protect\hypertarget{searching-for-a-design}{%
Searching for a design}}

\scriptsize

\begin{Shaded}
\begin{Highlighting}[]
\NormalTok{accrual.period =}\StringTok{ }\KeywordTok{c}\NormalTok{(}\DecValTok{3}\NormalTok{)}
\NormalTok{followup.period =}\StringTok{ }\KeywordTok{c}\NormalTok{(}\DecValTok{1}\NormalTok{,}\DecValTok{2}\NormalTok{,}\DecValTok{3}\NormalTok{,}\DecValTok{4}\NormalTok{)}
\NormalTok{num.accrual.periods =}\StringTok{ }\KeywordTok{length}\NormalTok{(accrual.period)}
\NormalTok{num.followup.periods =}\StringTok{ }\KeywordTok{length}\NormalTok{(followup.period)}
\NormalTok{perc.reduct =}\StringTok{ }\DecValTok{28}
\NormalTok{p =}\StringTok{ }\KeywordTok{matrix}\NormalTok{(}\DecValTok{0}\NormalTok{,}\DataTypeTok{nrow =}\NormalTok{ num.accrual.periods, }\DataTypeTok{ncol =}\NormalTok{ num.followup.periods)}
\ControlFlowTok{for}\NormalTok{(jj }\ControlFlowTok{in} \DecValTok{1}\OperatorTok{:}\NormalTok{num.followup.periods)\{}
\NormalTok{  p[}\DecValTok{1}\NormalTok{,jj]=}\StringTok{ }\KeywordTok{cpower}\NormalTok{(}\DataTypeTok{tref =} \FloatTok{0.75}\NormalTok{, }\DataTypeTok{n =} \DecValTok{150}\NormalTok{, }\DataTypeTok{mc =} \FloatTok{.50}\NormalTok{, }\DataTypeTok{r =}\NormalTok{ perc.reduct, }
       \DataTypeTok{accrual =}\NormalTok{ accrual.period, }\DataTypeTok{tmin =}\NormalTok{ followup.period[jj], }
       \DataTypeTok{noncomp.c =} \DecValTok{0}\NormalTok{, }\DataTypeTok{noncomp.i =} \DecValTok{0}\NormalTok{, }
       \DataTypeTok{alpha =} \FloatTok{0.05}\NormalTok{, }\DataTypeTok{pr =} \OtherTok{FALSE}\NormalTok{)}
\NormalTok{  \}}
\NormalTok{p}
\end{Highlighting}
\end{Shaded}

\begin{verbatim}
##         [,1]    [,2]    [,3]    [,4]
## [1,] 0.67402 0.72517 0.74768 0.75829
\end{verbatim}

\normalsize

A more thorough search\ldots

\end{frame}

\begin{frame}[fragile]

\scriptsize

\begin{Shaded}
\begin{Highlighting}[]
\NormalTok{accrual.period =}\StringTok{ }\KeywordTok{c}\NormalTok{(}\DecValTok{3}\NormalTok{,}\DecValTok{5}\NormalTok{,}\DecValTok{7}\NormalTok{,}\DecValTok{9}\NormalTok{)}
\NormalTok{enrollment.total =}\StringTok{ }\DecValTok{50}\OperatorTok{*}\NormalTok{accrual.period}
\NormalTok{followup.period =}\StringTok{ }\KeywordTok{c}\NormalTok{(}\DecValTok{1}\NormalTok{,}\DecValTok{3}\NormalTok{,}\DecValTok{5}\NormalTok{,}\DecValTok{7}\NormalTok{)}
\NormalTok{num.accrual.periods =}\StringTok{ }\KeywordTok{length}\NormalTok{(accrual.period)}
\NormalTok{num.followup.periods =}\StringTok{ }\KeywordTok{length}\NormalTok{(followup.period)}
\NormalTok{perc.reduct =}\StringTok{ }\DecValTok{28}
\NormalTok{p =}\StringTok{ }\KeywordTok{matrix}\NormalTok{(}\DecValTok{0}\NormalTok{,}\DataTypeTok{nrow =}\NormalTok{ num.accrual.periods, }\DataTypeTok{ncol =}\NormalTok{ num.followup.periods)}
\ControlFlowTok{for}\NormalTok{(ii }\ControlFlowTok{in} \DecValTok{1}\OperatorTok{:}\NormalTok{num.accrual.periods)\{}
   \ControlFlowTok{for}\NormalTok{(jj }\ControlFlowTok{in} \DecValTok{1}\OperatorTok{:}\NormalTok{num.followup.periods)\{}
\NormalTok{     p[ii,jj]=}\StringTok{ }\KeywordTok{cpower}\NormalTok{(}\DataTypeTok{tref =} \FloatTok{0.75}\NormalTok{, }\DataTypeTok{n =}\NormalTok{ enrollment.total[ii], }
                      \DataTypeTok{mc =} \FloatTok{.50}\NormalTok{, }\DataTypeTok{r =}\NormalTok{ perc.reduct, }
                      \DataTypeTok{accrual =}\NormalTok{ accrual.period[ii], }
                      \DataTypeTok{tmin =}\NormalTok{ followup.period[jj], }
                      \DataTypeTok{noncomp.c =} \DecValTok{0}\NormalTok{, }\DataTypeTok{noncomp.i =} \DecValTok{0}\NormalTok{, }
                      \DataTypeTok{alpha =} \FloatTok{0.05}\NormalTok{, }\DataTypeTok{pr =} \OtherTok{FALSE}\NormalTok{)}
\NormalTok{   \}}
\NormalTok{\}}
\NormalTok{power.table =}\StringTok{ }\KeywordTok{matrix}\NormalTok{(}\DecValTok{0}\NormalTok{, }\DataTypeTok{nrow =}\NormalTok{ num.accrual.periods }\OperatorTok{+}\StringTok{ }\DecValTok{1}\NormalTok{,}
                     \DataTypeTok{ncol =}\NormalTok{ num.followup.periods }\OperatorTok{+}\StringTok{ }\DecValTok{1}\NormalTok{)}
\NormalTok{power.table[}\DecValTok{1}\NormalTok{,] =}\StringTok{ }\KeywordTok{c}\NormalTok{(}\DecValTok{0}\NormalTok{, followup.period)}
\NormalTok{power.table[,}\DecValTok{1}\NormalTok{] =}\StringTok{ }\KeywordTok{c}\NormalTok{(}\DecValTok{0}\NormalTok{, accrual.period)}
\ControlFlowTok{for}\NormalTok{(ii }\ControlFlowTok{in} \DecValTok{1}\OperatorTok{:}\NormalTok{num.accrual.periods)\{}
  \ControlFlowTok{for}\NormalTok{(jj }\ControlFlowTok{in} \DecValTok{1}\OperatorTok{:}\NormalTok{num.followup.periods)\{}
\NormalTok{  power.table[ii }\OperatorTok{+}\StringTok{ }\DecValTok{1}\NormalTok{, jj }\OperatorTok{+}\StringTok{ }\DecValTok{1}\NormalTok{] =}\StringTok{ }\NormalTok{p[ii,jj]}
\NormalTok{  \}}
\NormalTok{\}}

\NormalTok{power.table }
\end{Highlighting}
\end{Shaded}

\end{frame}

\begin{frame}[fragile]

\scriptsize

\begin{Shaded}
\begin{Highlighting}[]
\KeywordTok{rownames}\NormalTok{(power.table) <-}\StringTok{ }\KeywordTok{c}\NormalTok{(}\StringTok{"yrs follow-up"}\NormalTok{, }\KeywordTok{rep}\NormalTok{(}\StringTok{"-"}\NormalTok{, }\DecValTok{4}\NormalTok{))}
\KeywordTok{colnames}\NormalTok{(power.table) <-}\StringTok{ }\KeywordTok{c}\NormalTok{(}\StringTok{"yrs enrollment"}\NormalTok{, }\KeywordTok{rep}\NormalTok{(}\StringTok{"-"}\NormalTok{, }\DecValTok{4}\NormalTok{))}
\KeywordTok{round}\NormalTok{(power.table, }\DataTypeTok{digits =} \DecValTok{3}\NormalTok{)}
\end{Highlighting}
\end{Shaded}

\begin{verbatim}
##               yrs enrollment     -     -     -     -
## yrs follow-up              0 1.000 3.000 5.000 7.000
## -                          3 0.674 0.748 0.764 0.768
## -                          5 0.900 0.928 0.934 0.935
## -                          7 0.975 0.982 0.984 0.984
## -                          9 0.994 0.996 0.997 0.997
\end{verbatim}

\normalsize

Studies often need to be larger than anticipated.

Take a closer look at the design with 5 years of enrollment and 3 years
of follow-up.

\end{frame}

\begin{frame}[fragile]

\scriptsize

\begin{verbatim}
## 
## Accrual duration: 5 y  Minimum follow-up: 3 y
## 
## Total sample size: 250 
## 
## Alpha= 0.05 
## 
## 0.75-year Mortalities
##      Control Intervention 
##         0.50         0.36 
## 
## Hazard Rates
##      Control Intervention 
##      0.92420      0.59505 
## 
## Probabilities of an Event During Study
##      Control Intervention 
##      0.98661      0.94649 
## 
## Expected Number of Events
##      Control Intervention 
##        123.3        118.3 
## 
## Hazard ratio: 0.64386 
## Standard deviation of log hazard ratio: 0.12869
\end{verbatim}

\begin{verbatim}
##   Power 
## 0.92803
\end{verbatim}

\end{frame}

\begin{frame}{%
\protect\hypertarget{what-is-the-effect-of-non-compliance}{%
What is the effect of non-compliance?}}

ITT: analyze according to assigned treatment, not treatment received.
\medskip

Main justification:

\begin{itemize}
\item
  \(p\)-values are calculated assuming no treatment difference (the null
  hypothesis)
\item
  Under that assumption, assigned treatment does not affect outcome
\item
  \(p\)-values will be correct (valid) when comparing the two groups
  according to treatment assignment
\end{itemize}

Example may help make this clear.

\end{frame}

\begin{frame}{%
\protect\hypertarget{simple-trial-no-difference-non-random-crossover}{%
Simple trial, no difference, non-random crossover}}

Suppose two treatments (\(A\) and \(B\)) are equally effective.

100 participants randomized to each treatment.

ITT table:

\begin{longtable}[]{@{}ccc@{}}
\toprule
\begin{minipage}[b]{0.15\columnwidth}\centering
Response\strut
\end{minipage} & \begin{minipage}[b]{0.22\columnwidth}\centering
Treatment \(A\)\strut
\end{minipage} & \begin{minipage}[b]{0.25\columnwidth}\centering
Treatment \(B\)\strut
\end{minipage}\tabularnewline
\midrule
\endhead
\begin{minipage}[t]{0.15\columnwidth}\centering
Success\strut
\end{minipage} & \begin{minipage}[t]{0.22\columnwidth}\centering
40\strut
\end{minipage} & \begin{minipage}[t]{0.25\columnwidth}\centering
40\strut
\end{minipage}\tabularnewline
\begin{minipage}[t]{0.15\columnwidth}\centering
Failure\strut
\end{minipage} & \begin{minipage}[t]{0.22\columnwidth}\centering
60\strut
\end{minipage} & \begin{minipage}[t]{0.25\columnwidth}\centering
60\strut
\end{minipage}\tabularnewline
\bottomrule
\end{longtable}

Now assume, after randomization:

\begin{itemize}
\item
  10 participants with good prognosis (future responders) switch from
  \(A\) to \(B\)
\item
  10 participants with bad prognosis (future non-responders) switch from
  \(B\) to \(A\)
\end{itemize}

\end{frame}

\begin{frame}{%
\protect\hypertarget{simple-trial-no-difference-non-random-crossover-1}{%
Simple trial, no difference, non-random crossover\ldots}}

Two treatments still equally effective.

Table for the as-treated groups:

\begin{longtable}[]{@{}ccc@{}}
\toprule
\begin{minipage}[b]{0.15\columnwidth}\centering
Response\strut
\end{minipage} & \begin{minipage}[b]{0.22\columnwidth}\centering
Treatment \(A\)\strut
\end{minipage} & \begin{minipage}[b]{0.25\columnwidth}\centering
Treatment \(B\)\strut
\end{minipage}\tabularnewline
\midrule
\endhead
\begin{minipage}[t]{0.15\columnwidth}\centering
Success\strut
\end{minipage} & \begin{minipage}[t]{0.22\columnwidth}\centering
30\strut
\end{minipage} & \begin{minipage}[t]{0.25\columnwidth}\centering
50\strut
\end{minipage}\tabularnewline
\begin{minipage}[t]{0.15\columnwidth}\centering
Failure\strut
\end{minipage} & \begin{minipage}[t]{0.22\columnwidth}\centering
70\strut
\end{minipage} & \begin{minipage}[t]{0.25\columnwidth}\centering
50\strut
\end{minipage}\tabularnewline
\bottomrule
\end{longtable}

An as-treated analysis would imply that \(B\) is more effective than
\(A\).

\end{frame}

\begin{frame}{%
\protect\hypertarget{simple-trial-difference-random-crossover}{%
Simple trial, difference, random crossover}}

ITT can be biased when there is a real treatment effect (random
crossovers).

Suppose \(B\) is more effective than \(A\), so for 100 in each group:

\begin{longtable}[]{@{}ccc@{}}
\toprule
\begin{minipage}[b]{0.15\columnwidth}\centering
Response\strut
\end{minipage} & \begin{minipage}[b]{0.22\columnwidth}\centering
Treatment \(A\)\strut
\end{minipage} & \begin{minipage}[b]{0.25\columnwidth}\centering
Treatment \(B\)\strut
\end{minipage}\tabularnewline
\midrule
\endhead
\begin{minipage}[t]{0.15\columnwidth}\centering
Success\strut
\end{minipage} & \begin{minipage}[t]{0.22\columnwidth}\centering
30\strut
\end{minipage} & \begin{minipage}[t]{0.25\columnwidth}\centering
50\strut
\end{minipage}\tabularnewline
\begin{minipage}[t]{0.15\columnwidth}\centering
Failure\strut
\end{minipage} & \begin{minipage}[t]{0.22\columnwidth}\centering
70\strut
\end{minipage} & \begin{minipage}[t]{0.25\columnwidth}\centering
50\strut
\end{minipage}\tabularnewline
\bottomrule
\end{longtable}

Assume 10 randomly chosen participants from each group switch treatments
after randomization but before starting treatment.

\end{frame}

\begin{frame}{%
\protect\hypertarget{table-with-only-patients-who-do-not-switch}{%
Table with only patients who do not switch}}

\begin{longtable}[]{@{}ccc@{}}
\toprule
\begin{minipage}[b]{0.15\columnwidth}\centering
Response\strut
\end{minipage} & \begin{minipage}[b]{0.22\columnwidth}\centering
Treatment \(A\)\strut
\end{minipage} & \begin{minipage}[b]{0.25\columnwidth}\centering
Treatment \(B\)\strut
\end{minipage}\tabularnewline
\midrule
\endhead
\begin{minipage}[t]{0.15\columnwidth}\centering
Success\strut
\end{minipage} & \begin{minipage}[t]{0.22\columnwidth}\centering
27\strut
\end{minipage} & \begin{minipage}[t]{0.25\columnwidth}\centering
45\strut
\end{minipage}\tabularnewline
\begin{minipage}[t]{0.15\columnwidth}\centering
Failure\strut
\end{minipage} & \begin{minipage}[t]{0.22\columnwidth}\centering
63\strut
\end{minipage} & \begin{minipage}[t]{0.25\columnwidth}\centering
45\strut
\end{minipage}\tabularnewline
\bottomrule
\end{longtable}

Attrition did not change measured success rates, but\ldots{}

\begin{itemize}
\tightlist
\item
  does reduce the effective sample size
\end{itemize}

What happens when `switchers’ are put back in?

\begin{itemize}
\item
  10 \(A\) \(\rightarrow\) \(B\), 5 respond, 5 do not
\item
  10 \(B\) \(\rightarrow\) \(A\), 3 respond, 7 do not
\end{itemize}

\end{frame}

\begin{frame}{%
\protect\hypertarget{itt-table-with-assigned-treatment-real-response}{%
ITT table with assigned treatment, real response}}

\(A\) gets 5 responders (who received B)

\(B\) gets 3 responders (who received A)

\begin{longtable}[]{@{}ccc@{}}
\toprule
\begin{minipage}[b]{0.15\columnwidth}\centering
Response\strut
\end{minipage} & \begin{minipage}[b]{0.22\columnwidth}\centering
Treatment \(A\)\strut
\end{minipage} & \begin{minipage}[b]{0.25\columnwidth}\centering
Treatment \(B\)\strut
\end{minipage}\tabularnewline
\midrule
\endhead
\begin{minipage}[t]{0.15\columnwidth}\centering
Success\strut
\end{minipage} & \begin{minipage}[t]{0.22\columnwidth}\centering
32\strut
\end{minipage} & \begin{minipage}[t]{0.25\columnwidth}\centering
48\strut
\end{minipage}\tabularnewline
\begin{minipage}[t]{0.15\columnwidth}\centering
Failure\strut
\end{minipage} & \begin{minipage}[t]{0.22\columnwidth}\centering
68\strut
\end{minipage} & \begin{minipage}[t]{0.25\columnwidth}\centering
52\strut
\end{minipage}\tabularnewline
\bottomrule
\end{longtable}

Apparent success rate:

\begin{itemize}
\item
  \(A\): 32\% after crossover vs.~30\% before crossover
\item
  \(B\): 48\% after crossover vs 50\% before crossover
\end{itemize}

Response proportions have moved closer together.

Non-random attrition can also cause bias in the analysis because of
missing data.

\end{frame}

\begin{frame}{%
\protect\hypertarget{accounting-for-non-compliance}{%
Accounting for non-compliance}}

If some patients do not take their assigned treatments, power of the
study will decrease. This issue has two sides:

\emph{Drop-outs} (\(d_e\))

\begin{itemize}
\item
  Patients who cannot tolerate the medication stop taking it.
\item
  Their hazard rate would become the same as the placebo group (if
  included in study) at that point.
\end{itemize}

\emph{Drop-ins} (\(d_c\))

\begin{itemize}
\tightlist
\item
  Patients assigned to less effective therapy may not get symptom relief
  and seek other therapy, or request to cross over.
\end{itemize}

A conservative remedy: adjust \(P_e\) and \(P_c\) as follows:

\begin{align*}
P_e^{*} &= P_e (1-d_e) + P_c d_e\\[1ex]
P_c^{*} &=  P_c (1-d_c) + P_e d_c
\end{align*}

\end{frame}

\begin{frame}{%
\protect\hypertarget{esophageal-cancer-again}{%
Esophageal cancer, again}}

\begin{itemize}
\item
  Treatments for cancer are often toxic, and patients have difficulty
  adhering to regimens.
\item
  \texttt{cpower()} can incoporate non-adherence through the parameters
  \texttt{noncomp.c, noncomp.i}.
\item
  What happens to power with 20\% non-adherence on the intervention arm?
\end{itemize}

\end{frame}

\begin{frame}[fragile]

\scriptsize

\begin{Shaded}
\begin{Highlighting}[]
\KeywordTok{library}\NormalTok{(Hmisc)}
\KeywordTok{cpower}\NormalTok{(}\DataTypeTok{tref =} \FloatTok{0.75}\NormalTok{, }\DataTypeTok{n =} \DecValTok{50}\OperatorTok{*}\DecValTok{5}\NormalTok{, }\DataTypeTok{mc =} \FloatTok{0.5}\NormalTok{, }
       \DataTypeTok{r =} \DecValTok{28}\NormalTok{, }\DataTypeTok{accrual=} \DecValTok{3}\NormalTok{, }\DataTypeTok{tmin =} \DecValTok{3}\NormalTok{, }\DataTypeTok{noncomp.c=}\DecValTok{0}\NormalTok{, }\DataTypeTok{noncomp.i=}\DecValTok{20}\NormalTok{, }
       \DataTypeTok{alpha=}\FloatTok{0.05}\NormalTok{, }\DataTypeTok{pr=}\OtherTok{FALSE}\NormalTok{)}
\end{Highlighting}
\end{Shaded}

\begin{verbatim}
##   Power 
## 0.77438
\end{verbatim}

\end{frame}

\begin{frame}{%
\protect\hypertarget{what-about-attrition}{%
What about attrition}}

\emph{Attrition} happens when participants leave a study; treatment and
outcome assessments cannot be made.

Attrition is also called \emph{loss to follow-up} and is common in long
term follow-up for studies of chronic diseases.

\texttt{cpower()} cannot adjust for attrition, but \texttt{nSurv()} in
the package \texttt{gsDesign} can.

A lab exercise uses \texttt{nSurv()} to reproduce the design of the
SPRINT trial summarized at the beginning of this unit.

\end{frame}

\hypertarget{derivations}{%
\section{Derivations}\label{derivations}}

\begin{frame}{%
\protect\hypertarget{heuristic-proof-of-the-distribution-of-the-log-rank}{%
Heuristic proof of the distribution of the log-rank}}

\begin{align*} 
 E(d_{1j} | d_{1j}, d_{0j}, r_{1j}, r_{0j}) 
    &= P( d_{1j} = 1 | d_{j}=1, r_{1j}, r_{0j}) \\[1ex]
    &=  \frac  {r_{1j} \lambda_0 \theta}
    {r_{1j} \lambda_0 \theta + r_{0j} \lambda_0}\\[1ex]
     &=  \frac  {r_{1j}  \theta} {r_{1j}  
\theta + r_{0j} }\\[1ex]
     &=  \frac  {r_{1j}} {r_{1j} + r_{0j} } + \log(\theta) 
      \left[\frac  {r_{1j}r_{0j}} {(r_{1j} + r_{0j})^2}\right]  \\
\end{align*}

But \(e_j = r_{1j}/(r_{1j} + r_{0j})\), so:
\[    E(d_{1j} | d_{1j}, d_{0j}, r_{1j}, r_{0j}) - e_j =
 \log(\theta) \left[\frac  {r_{1j}r_{0j}} {(r_{1j} + r_{0j})^2 }\right] \]

\end{frame}

\begin{frame}{%
\protect\hypertarget{heuristic-proof-of-the-distribution}{%
Heuristic proof of the distribution\ldots}}

If \(n_0=n_1\), then near H\(_0\), \(r_{1j} \approx r_{0j}\), hence,
\[    E(d_{1j} | d_{1j}, d_{0j}, r_{1j}, r_{0j}) - e_j = \log(\theta) /4 \]

Similarly, with no ties, \[ v_j = r_{1j} r_{0j}/r_j^2 \approx 1/4 \]

\end{frame}

\begin{frame}{%
\protect\hypertarget{alternative-derivation-using-the-partial-likelihood}{%
Alternative derivation using the partial likelihood}}

\small

The partial likelihood is: \begin{align*}
  l(\bbeta)   &=  \log \left[ \prod_{j=1}^{n}\left(    \frac
      {   e^{\bbeta \bZ_j } }   {\sum_{\ell \in {\cal R}(\tau_j)}
e^{\bbeta \bZ_{\ell}} }\right)^{\delta_j} \right]  \\[1ex]
&=\sum_{j=1}^{n} \delta_j  \left[ \bbeta \bZ_j  -
\log \left(\sum_{\ell \in {\cal R}(\tau_j)} e^{\bbeta \bZ_{\ell}}\right) 
\right]
\end{align*}

The partial derivative of log-likelihood (the score statistic) is:
\begin{align*}
U(\bbeta)   &=  \frac{\partial}{\partial \bbeta} \ell(\bbeta) \\
&=  \sum_{j=1}^{n} \delta_j  \left[ \bZ_j  -
\frac{ \sum_{\ell \in {\cal R}(\tau_j)} \bZ_{\ell} ~e^{\bbeta \bZ_{\ell}} }
{ \sum_{\ell \in {\cal R}(\tau_j)}   e^{\bbeta  \bZ_{\ell}} } \right]
\end{align*}

\end{frame}

\begin{frame}{%
\protect\hypertarget{partial-likelihood-derivation}{%
Partial likelihood derivation \ldots}}

The negative second partial derivative of the log-likelihood) is:
\footnotesize \[
-\frac{\partial^2}{\partial \bbeta^2} \ell(\bbeta) =
 \sum_{j=1}^{n}   \delta_j \left[     
\frac{ \sum_{\ell \in {\cal R}(\tau_j)} e^{\bbeta \bZ_{\ell}}
\sum_{\ell \in {\cal R}(\tau_j)} \bZ_{\ell} e^{\bbeta \bZ_{\ell}}-
(\sum_{\ell \in {\cal R}(\tau_j)} \bZ_{\ell} e^{\bbeta \bZ_{\ell}})^2}
{ \sum_{\ell \in {\cal R}(\tau_j)} e^{\bbeta  \bZ_{\ell}} } \right]
\]

\normalsize

The logrank statistic (with no ties) is equivalent to the score
statistic for testing \(\bbeta=0\): \[
Z_{LR} = \frac{U(0)}{\sqrt{I(0)}}
\] By a Taylor series expansion:
\[U(0) \cong U(\beta)- \beta \frac{\partial U}{\partial \beta}(0)\]
\[E[U(0)]\cong \beta d /4 \hspace{0.3in} \text{and ~} I(0)\cong d/4\]

\end{frame}

\end{document}
